\setlength{\parskip}{0em}
I declare that the work described in this thesis was carried out by me and has not been presented in any previous application for a degree at the University of Glasgow or any other institution. The chapters in this thesis and the contributions of others to them are described below.
\paragraph{Chapter \ref{chapter:Gravitational_wave_interferometers}}
This chapter contains introductory information pertaining to ground-based gravitational wave interferometry and the LIGO detectors.
\paragraph{Chapter \ref{chapter:TnT_lab_work}} 
In this chapter, experiments to do with the prototype of the laser that will be used for LIGO's fourth observation run are described. My contribution to this work was in the setup of the optical table and the taking of measurements. These tasks were completed with help from Nina Bode, Matthew Heintze and Michael Fyffe. Rick Savage and Maik Frede contributed to discussions of how this should be done. The diagnostic breadboard was supplied by the AEI. The design of the prototype was undertaken by others. The outcome of this project was published in \cite{Bode2020a}. The photographs of the laser were taken by Michael Fyffe. This work was carried out at the LIGO Livingston Observatory as part of the LSC fellows programme and was funded by the STFC long-term attachment scheme.
\paragraph{Chapter \ref{chapter:simulation_finesse}}
In this chapter, simulations of the LIGO Livingston interferometer were performed to investigate the effect of mismatch between the arm cavities and the output mode cleaners. This work was aided by discussions with Ken Strain, Stephen Webster, Bryan Barr and Carl Blair. 
\paragraph{Chapter \ref{chapter:BHD_for_ligo_1}} 
In this chapter, a review is given for the motivations and phase noise requirement of the balanced homodyne detection scheme that will be installed as part of the LIGO A+ detectors.
\paragraph{Chapter \ref{chapter:BHD_for_ligo_2}} 
In this chapter, a method for showing the effectiveness of the active optics for adjusting the mode matching between the interferometer and the output mode cleaners for the LIGO A+ detectors is described. With this method, I showed the range of mode mismatches that the active optics in HAM6 will be able to correct. The design of the optical layout was undertaken by me, Ken Strain, Mark Barton, Russell Jones and Stephen Webster. As part of this, I wrote an algorithm for finding mode matching solutions; Stephen Webster provided useful input into this.
\paragraph{Chapter \ref{chapter:2um_photoiode_chapter}}
The experiments reported on in this chapter were done under the supervision of Ken Strain and Bryan Barr. Sebastian Steinlechner supplied the IG24x500S4i and IG26x500S4i photodiodes. I would like to thank Sean Leavey and Andrew Long for their helpful discussions about photodiodes.
\paragraph{Chapter \ref{chapter:in_air_MZ}}
The in-air apparatus was constructed with help from Russell Jones and Stephen Craig. The laser, send/receive circuits and CDS infrastructure were repurposed from the MOSES experiment at the Glasgow 10\,m prototype interferometer worked on by Bryan Barr, Neil Gordon, Jennifer Wright, Borja Sorazu and Ken Strain. The high voltage power supply used to drive the PZT in this experiment was designed by Andrew Spencer and Steven O'Shea. The in-vacuum experiment was performed with repurposed apparatus left over from a previous experiment run by Stefan Hild. The auxiliary suspensions were designed by Jan-Simon Hennig and Russell Jones. The balanced homodyne suspension and coil drivers were designed by Sebastian Steinlechner.
\setlength{\parskip}{2em}
\paragraph{Other Credits}
Many of the figures in this thesis were created with Matplotlib~\cite{Caswell2021}. Diagrams of optical layouts were often created with an open source collection of graphics produced by Alexander Franzen~\cite{ComponentLibrary}. The work in this thesis was facilitated by the University of Glasgow and was funded by the STFC (award number: 1947199).